\documentclass[10pt, twocolumn, a4paper]{article}

% Codificação e idioma
\usepackage[utf8]{inputenc}
\usepackage[T1]{fontenc}
\usepackage[brazil]{babel}

% Formatação ABNT (margens)
\usepackage[
  top=3cm,
  bottom=2cm,
  left=3cm,
  right=2cm
]{geometry}

% Espaçamento entre linhas (ABNT: 1,5 para texto, mas twocolumn usa simples)
\usepackage{setspace}
\onehalfspacing

% Fontes e tipografia
\usepackage{times}
\usepackage{microtype}

% Outros pacotes
\usepackage{indentfirst}
\usepackage{url}
\usepackage{hyperref}

% Referências ABNT (deve vir após hyperref)
\usepackage[alf]{abntex2cite}

% Configuração de hyperlinks
\hypersetup{
  colorlinks=true,
  linkcolor=black,
  citecolor=black,
  urlcolor=black
}

% Título, autor e data
\title{\textbf{Confiabilidade de sistemas de software distribuídos a partir de sinais OpenTelemetry}}
\author{Josoe Santos}
\date{2026}

\begin{document}

\maketitle

\begin{abstract}
Este trabalho propõe o estudo da confiabilidade de sistemas de software distribuídos utilizando sinais coletados pelo padrão aberto OpenTelemetry. A iniciativa visa democratizar práticas de \textit{Software Reliability Engineering} (SRE) para médias e pequenas empresas brasileiras, reduzindo a dependência de fornecedores proprietários de monitoramento.
\end{abstract}

\section{Introdução}

O tema de confiabilidade de software é relativamente recente, mas amplamente discutido. A ubiquidade de sistemas de software em diversas áreas da sociedade tornou-os complexos e mais críticos. Nas décadas mais recentes, com os avanços no campo da internet e sistemas de software dependentes dela, tornou-se prática comum separar diferentes partes complexas de aplicações em servidores conectados por redes \cite{fowler2002patterns}, isto é o que denominamos sistemas distribuídos. A partir dessa definição, cada servidor em um sistema distribuído é tratado como um componente que pode falhar ou se comportar fora do esperado, similarmente a componentes de outros tipos de artefatos de engenharia.

A prática que estuda a confiabilidade de sistemas web é conhecida como \textit{Software Reliability Engineering} (SRE) \cite{google_sre} e é prática comum no Brasil desde a metade da última década. Contudo, esse campo ainda não é consistentemente aplicado em empresas brasileiras, porque o tema demanda alto investimento e planejamento estratégico de longo prazo, tornando-se inviável para médias e pequenas empresas. Os custos relacionados a tempo, pessoas e fornecedores são altos demais para uma adoção da prática que realmente impacte positivamente a confiabilidade de sistemas novos.

Normalmente, um dos fornecedores mais caros da SRE é o fornecedor de monitoramento. Essas empresas têm como maior produto uma suíte de bibliotecas e ferramentas de nuvem especializadas em coletar e processar dados de monitoramento de sistemas distribuídos. Contudo, o número relativamente pequeno de fornecedores competitivos leva a preços altos e ao efeito de \textit{vendor lock-in}, que é a incapacitação de uma organização de substituir o fornecedor por ser complexo ou caro demais. Essa situação irrita o mercado de software há anos e, para tanto, foi criado o OpenTelemetry \cite{otel_site}, uma iniciativa de código aberto que tem como finalidade padronizar e simplificar a adoção de software de monitoramento.

\section{Descrição do sistema}

A proposta de trabalho é estruturada em três etapas incrementais:

\begin{enumerate}
  \item Iniciar com um servidor simples e demonstrar a \textit{stack} LGTM (Loki, Grafana, Tempo, Mimir) com base nos dados de um sistema;
  \item Adicionar um segundo servidor e mostrar os princípios se repetindo;
  \item Aplicar o padrão para $N$ servidores interconectados com ou sem redundância. Nessa fase, podem ser utilizados algoritmos de grafos e as fórmulas de probabilidade combinada (independente e dependente).
\end{enumerate}

\section{Revisão de Literatura}

A revisão de literatura contemplará as seguintes fontes:

\begin{itemize}
  \item Livros de SRE do Google \cite{google_sre};
  \item Artigos da IEEE sobre confiabilidade de sistemas distribuídos;
  \item Artigos da Thoughtworks (martinfowler.com) \cite{fowler2002patterns}.
\end{itemize}

\section{Metodologia}

A metodologia proposta compreende as seguintes etapas:

\begin{enumerate}
  \item Coleta de dados por 30 dias;
  \item Padronização de sinais pelo OpenTelemetry;
  \item Desenvolvimento de ferramentas de baixo custo para uso no Grafana.
\end{enumerate}

\section{Resultados Esperados}

Espera-se como resultado deste trabalho:

\begin{itemize}
  \item Visão mais realista da saúde do sistema — os resultados devem estar em consonância com a opinião dos usuários;
  \item Melhoria da visibilidade de problemas para usuários do Grafana.
\end{itemize}

\bibliography{references}

\end{document}
